\section{Related Work}
\label{sec:related_work}

\subsection{Long-Term Agent Memory}
\label{subsec:rw-agent-memory}

Long-term agent memory systems externalize past interaction into persistent records, summaries, skills, rules, and state.
Our work is complementary to this line: it does not replace compression or storage mechanisms, but asks how their outputs are accessed and exposed at runtime.

\subsection{Retrieval and Structured Memory}
\label{subsec:rw-retrieval}

Retrieval-augmented generation and graph-based memory systems provide access methods beyond full-context prompting.
Systems such as hierarchical graph memory and direct corpus interaction suggest that access language matters: top-$k$ similarity search, graph browsing, shell-like corpus operations, and workspace interaction expose different information under different budgets.
Our framework treats these as different primitive sets rather than as interchangeable retrieval variants \cite{tang2026mnemis}.

\subsection{Benchmarks for Memory Use}
\label{subsec:rw-benchmarks}

Personalized memory QA benchmarks test whether agents can answer referential questions grounded in a user's past.
State-invalidation benchmarks test whether agents can revise old beliefs in light of later observations.
We use these benchmarks not only for task scores, but as controlled environments for decomposing where memory-use failures occur.

\subsection{Evidence, Provenance, and Traceability}
\label{subsec:rw-evidence}

Work on evidence grounding, provenance, execution tracing, and failure attribution is closely related to our notion of evidence views.
The distinction is that we focus on what structure must be visible to the downstream decision model, rather than only whether the system can recover a post-hoc explanation.
